% Part: counterfactuals
% Chapter: minimal-change-semantics
% Section: transitivity

\documentclass[../../../include/open-logic-section]{subfiles}

\begin{document}

\olfileid[pt]{cnt}{min}{cpo}

\olsection{Contraposição}

Condicionais materiais e estritos equivalem às suas
contrapositivas. Contrafactuais não. Eis um exemplo de
Kratzer:
\begin{quote}
  Se Goethe não tivesse morrido em 1832, ele (ainda) estaria morto agora.

  Se Goethe não estivesse morto agora, ele teria morrido em 1832.
\end{quote}
A primeira sentença é verdadeira: seres humanos não vivem centenas de anos. A
segunda é claramente falsa: se Goethe não estivesse morto agora, ainda estaria
vivo e, portanto, não poderia ter morrido em 1832.

\begin{ex}\ollabel{ex:contraposition-counterex}
  A semântica de esferas invalida a contraposição, O.s.\@, temos $p
  \cif q \Entails/ \lnot q \cif \lnot p$. Interprete $p$ como ``Goethe
  não morreu em 1832'' e $q$ como ``Goethe está morto agora''. Podemos
  representar isso em um modelo $\mModel{M_1} = \tuple{W, O, V}$ com $W =
  \{w, w_1, w_2\}$, $O = \{\{w\}, \{w, w_1\}, \{w, w_1, w_2\}\}$,
  $V(p) = \{w_1, w_2\}$ e $V(q) = \{w, w_1\}$. Assim, $w$ é o mundo
  real em que Goethe morreu em 1832 e continua morto; $w_1$ é o
  mundo próximo em que Goethe morreu, digamos, em 1833 e continua morto;
  e $w_2$ é um mundo remoto em que Goethe ainda está vivo.
\begin{figure}
\begin{center}
\begin{tikzpicture}[scale=.6,layerwidth=1.5]\tiny
  \spheresystem{3}
  \begin{scope}
    \clip \propositionplot{0}{.8}{50}{4.5} -- (4.5,-4.5) -- (-4.5,-4.5) -- (-4.5,4.5) -- (4.5,4.5);
    \spherelayer{3}
  \end{scope}
  \propositionintersect{0}{2}{110}{5.5}
  \proposition{0}{.8}{50}{4.5}
  \path (0,0) node[world] {$w$};
  \spherepos{0}{2}{node[world] {$w_1$}}
  \spherepos{-50}{3}{node[world] {$w_2$}}
  \spherepos{28}{3}{node {$q$}}
  \spherepos{42}{3}{node {$\lnot q$}}
  \spherepos{-55}{4}{node {$p$}}
  \spherepos{-67}{4}{node {$\lnot p$}}
\end{tikzpicture}
\caption{Contraexemplo à contraposição}
\ollabel{fig:contraposition}
\end{center}
\end{figure}
Existe uma esfera que admite $p$, $S =
\{w, w_1\}$, e $p \lif q$ é verdadeiro em todos os mundos nela, portanto $\mSat{M}{p
  \cif q}[w]$. Contudo, a esfera que admite $\lnot q$, $\{w, w_1,
w_2\}$, contém um mundo, a saber,~$w_2$, em que $q$ é falso e $p$ é
verdadeiro, portanto $\mSat/{M}{\lnot q \lif \lnot p}[w_2]$.
\end{ex}

\end{document}
